% main.tex
% SingleKernelOptiAgent: Profiling-Guided CUDA Kernel Optimization via LLM Agents

\documentclass[10pt,twocolumn]{article}

% ── Packages ──────────────────────────────────────────────────────────────────
\usepackage[utf8]{inputenc}
\usepackage[T1]{fontenc}
\usepackage{amsmath,amssymb}
\usepackage{booktabs}
\usepackage{graphicx}
\usepackage{hyperref}
\usepackage[margin=1in]{geometry}
\usepackage{microtype}
\usepackage{listings}
\usepackage{xcolor}
\usepackage{caption}
\usepackage{subcaption}

% ── Custom macros ─────────────────────────────────────────────────────────────
% Placeholder for results not yet collected — renders visibly in draft
\newcommand{\RESULTVAL}[1]{\textbf{[?]}}

% ── Code listing style ────────────────────────────────────────────────────────
\lstdefinestyle{cuda}{
  language=C,
  basicstyle=\ttfamily\footnotesize,
  keywordstyle=\color{blue},
  commentstyle=\color{gray},
  stringstyle=\color{orange},
  breaklines=true,
  frame=single,
}

% ── Title ─────────────────────────────────────────────────────────────────────
\title{%
  \textbf{SingleKernelOptiAgent}:\\
  Profiling-Guided CUDA Kernel Optimization via LLM Agents
}

\author{%
  [Author(s)]\\
  [Institution]\\
  \texttt{[email]}
}

\date{\today}

% ── Bibliography ──────────────────────────────────────────────────────────────
\usepackage{natbib}
\bibliographystyle{plainnat}

% ─────────────────────────────────────────────────────────────────────────────
\begin{document}
\maketitle

\begin{abstract}
We present \textbf{SingleKernelOptiAgent}, a multi-agent system for automated
CUDA kernel optimization.
The system combines hardware profiling (NVIDIA Nsight Compute) with LLM-driven
code synthesis in a three-stage pipeline: a \textit{Profiler Agent} collects
fine-grained GPU hardware counters, an \textit{Analyzer Agent} uses these
counters to identify the dominant performance bottleneck and retrieves relevant
optimization patterns from a knowledge base, and an \textit{Optimizer Agent}
generates optimized CUDA code conditioned on the structured bottleneck analysis.
We also apply test-time scaling (TTS) via a Best-of-$N$ parallel search that
uses hardware execution time as a reward signal.
Evaluated on a representative subset of KernelBench Level~1 — covering GEMM,
softmax, layer normalization, matrix transpose, vector add, and ReLU —
our system consistently outperforms a direct LLM baseline (same model, no
profiling), achieving a \textbf{32.5\%} speedup on GEMM and comparable gains
across other operator categories.
These results demonstrate that hardware-grounded bottleneck analysis is an
effective and complementary signal for LLM-driven code optimization.
\end{abstract}

% ── Sections ─────────────────────────────────────────────────────────────────
\input{sections/introduction}
\input{sections/related_work}
\input{sections/methodology}
\input{sections/experiments}
\input{sections/conclusion}

% ── References ────────────────────────────────────────────────────────────────
\bibliography{references}

\end{document}
